\chapter{Chapter06 — Random Small Bias}
%%% generated by the blueprint pipeline from FABL.Chapter06.Constructions.RandomSmallBias
\section{Overview}

Declarations of the Lean module \texttt{FABL.Chapter06.Constructions.RandomSmallBias}.

\section{Declarations}

\begin{definition}[Random small-bias constant]
\lean{FABL.randomSmallBiasConstant}
\label{FABL.randomSmallBiasConstant}
\leanok
The random small-bias constant is the natural number $8$.
\end{definition}

\begin{definition}[Random small-bias sample count]
\lean{FABL.randomSmallBiasSampleCount}
\label{FABL.randomSmallBiasSampleCount}
\leanok
\uses{FABL.randomSmallBiasConstant}
For a natural number $n$ and a real number $\varepsilon$, the sample count is the least
natural number greater than or equal to $8n/\varepsilon^2$.
\end{definition}

\begin{theorem}[Positivity of the sample count]
\lean{FABL.randomSmallBiasSampleCount_pos}
\label{FABL.randomSmallBiasSampleCount_pos}
\leanok
\uses{FABL.randomSmallBiasConstant, FABL.randomSmallBiasSampleCount}
\difficulty{3}
Let $n$ be a positive natural number and let $\varepsilon>0$. The least natural number
greater than or equal to $8n/\varepsilon^2$ is positive.
\end{theorem}

\begin{theorem}[Zero sample count at dimension zero]
\lean{FABL.randomSmallBiasSampleCount_zero}
\label{FABL.randomSmallBiasSampleCount_zero}
\leanok
\uses{FABL.randomSmallBiasSampleCount}
\difficulty{1}
For every real number $\varepsilon$, the random small-bias sample count at dimension
zero is $0$.
\end{theorem}

\begin{definition}[Random small-bias sample multiset]
\lean{FABL.randomSmallBiasMultiset}
\label{FABL.randomSmallBiasMultiset}
\leanok
\uses{FABL.F₂Cube}
Given natural numbers $n,m$ and a function assigning a binary cube of length $n$ to each
index in $\{0,\ldots,m-1\}$, form the multiset of those $m$ cube values, retaining
repetitions.
\end{definition}

\begin{theorem}[Cardinality of the sample multiset]
\lean{FABL.randomSmallBiasMultiset_card}
\label{FABL.randomSmallBiasMultiset_card}
\leanok
\uses{FABL.F₂Cube, FABL.randomSmallBiasMultiset}
\difficulty{1}
Let $n,m$ be natural numbers, and choose any list of $m$ vectors in $(\mathbb{F}_2)^n$.
The multiset formed by these samples, with repetitions counted, has cardinality $m$.
\end{theorem}

\begin{definition}[Small-bias sample]
\lean{FABL.IsSmallBiasedSample}
\statementsource{Analysis-of-Boolean-Functions-by-Ryan-ODonnell}{pdf-ffbe9c7b3306-p144-b007}
\label{FABL.IsSmallBiasedSample}
\leanok
\uses{FABL.F₂Cube, FABL.finiteUniformEmpiricalMean, FABL.vectorWalshCharacter}
Let $n,m$ be nonnegative integers, let $x_0,\ldots,x_{m-1}$ be vectors in the binary
cube of dimension $n$, and let $\varepsilon\in\mathbb{R}$. The sample is small-biased
with bias at most $\varepsilon$ if, for every nonzero binary vector $\gamma$ of length
$n$, the absolute value of the empirical mean of the Walsh character indexed by $\gamma$
on the sample is at most $\varepsilon$.
\end{definition}

\begin{theorem}[Small bias in dimension zero]
\lean{FABL.isSmallBiasedSample_zero}
\label{FABL.isSmallBiasedSample_zero}
\leanok
\uses{FABL.F₂Cube, FABL.IsSmallBiasedSample}
\difficulty{2}
For any natural number $m$, any sample of $m$ elements of the zero-dimensional vector
space over $\mathbb F_2$, and any real number $\varepsilon$, the sample is
$\varepsilon$-biased: every nonzero vector frequency has empirical character mean of
absolute value at most $\varepsilon$.
\end{theorem}

\begin{definition}[Empirical probability density]
\lean{FABL.randomSmallBiasDensity}
\statementsource{Analysis-of-Boolean-Functions-by-Ryan-ODonnell}{
  pdf-ffbe9c7b3306-p144-b007,
  pdf-ffbe9c7b3306-p153-b002
}
\label{FABL.randomSmallBiasDensity}
\leanok
\uses{FABL.F₂Cube, FABL.ProbabilityDensity, FABL.ProbabilityDensity.uniformPushforward}
Let $n,m$ be nonnegative integers with $m>0$, and let $x_0,\ldots,x_{m-1}$ be vectors in
the binary cube of dimension $n$. The associated probability density on that cube is the
density of the uniform distribution on the sample indices pushed forward by the map
sending each index to its sample vector.
\end{definition}

\begin{theorem}[Fourier coefficient of an empirical density]
\lean{FABL.vectorFourierCoeff_randomSmallBiasDensity}
\label{FABL.vectorFourierCoeff_randomSmallBiasDensity}
\leanok
\uses{FABL.F₂Cube, FABL.ProbabilityDensity.expectation, FABL.ProbabilityDensity.expectation_uniformPushforward, FABL.ProbabilityDensity.uniformPushforward, FABL.finiteUniformEmpiricalMean, FABL.randomSmallBiasDensity, FABL.vectorFourierCoeff, FABL.vectorFourierCoeff_eq_expect, FABL.vectorWalshCharacter}
\difficulty{3}
Let $m>0$, let $x_0,\ldots,x_{m-1}$ be a list of samples from the $n$-dimensional binary
cube, and let $\gamma$ be another point of that cube. The vector Fourier coefficient of
the density assigning mass according to these samples equals the empirical mean, over
the $m$ samples, of the Walsh character indexed by $\gamma$.
\end{theorem}

\begin{theorem}[Random small bias density is biased iff]
\lean{FABL.randomSmallBiasDensity_isBiased_iff}
\statementsource{Analysis-of-Boolean-Functions-by-Ryan-ODonnell}{pdf-ffbe9c7b3306-p144-b007}
\label{FABL.randomSmallBiasDensity_isBiased_iff}
\leanok
\uses{FABL.F₂Cube, FABL.IsSmallBiasedSample, FABL.ProbabilityDensity.IsBiased, FABL.randomSmallBiasDensity, FABL.vectorFourierCoeff_randomSmallBiasDensity}
\difficulty{3}
% dataset-pipeline: stub -- prose pending build_informal_statements.py
Statement pending regeneration.
\end{theorem}

\begin{definition}[Frequency bad set]
\lean{FABL.randomSmallBiasFrequencyBadSet}
\label{FABL.randomSmallBiasFrequencyBadSet}
\leanok
\uses{FABL.F₂Cube, FABL.finiteUniformEmpiricalMean, FABL.randomSmallBiasSampleCount, FABL.vectorWalshCharacter}
For a natural number $n$, real number $\varepsilon$, and binary vector $\gamma$ of
length $n$, the frequency bad set consists of sample sequences for which $\gamma$ is
nonzero and the absolute empirical mean of the Walsh character indexed by $\gamma$ is at
least $\varepsilon$.
\end{definition}

\begin{definition}[Small-bias failure set]
\lean{FABL.randomSmallBiasFailureSet}
\label{FABL.randomSmallBiasFailureSet}
\leanok
\uses{FABL.F₂Cube, FABL.IsSmallBiasedSample, FABL.randomSmallBiasSampleCount}
For a natural number $n$ and real number $\varepsilon$, the failure set consists of
sample sequences of binary cubes of length $n$ that are not $\varepsilon$-small-biased.
\end{definition}

\begin{theorem}[A bad-set frequency bound]
\lean{FABL.measure_randomSmallBiasFrequencyBadSet_le}
\label{FABL.measure_randomSmallBiasFrequencyBadSet_le}
\leanok
\uses{FABL.F₂Cube, FABL.expect_vectorWalshCharacter, FABL.measure_finiteUniformEmpiricalMean_sub_expect_ge_le, FABL.randomSmallBiasConstant, FABL.randomSmallBiasFrequencyBadSet, FABL.randomSmallBiasSampleCount, FABL.randomSmallBiasSampleCount_pos, FABL.uniformPMF, FABL.vectorWalshCharacter, FABL.vectorWalshCharacter_eq_neg_one_or_one}
\difficulty{5}
Let $n$ be a positive integer, let $\varepsilon>0$, and let $\gamma$ be a point of the
$n$-dimensional cube over $\mathbb F_2$. Choose $\lceil 8n/\varepsilon^2\rceil$ points
independently and uniformly from that cube. The probability that $\gamma\ne0$ and the
absolute empirical mean of the Walsh character indexed by $\gamma$ is at least
$\varepsilon$ is at most $2\exp(-4n)$.
\end{theorem}

\begin{theorem}[An exponential power bound]
\lean{FABL.two_pow_mul_two_exp_neg_four_mul_le_rpow_neg}
\label{FABL.two_pow_mul_two_exp_neg_four_mul_le_rpow_neg}
\leanok
\uses{FABL.two_pow_mul_two_exp_neg_two_mul_le_rpow_neg}
\difficulty{4}
For every integer $n\ge 1$, $2^n\,2e^{-4n}\le 2^{-n}$.
\end{theorem}

\begin{theorem}[Small bias of a random sample]
\lean{FABL.measure_randomSmallBiasSample_not_isBiased_le}
\label{FABL.measure_randomSmallBiasSample_not_isBiased_le}
\leanok
\uses{FABL.F₂Cube, FABL.IsSmallBiasedSample, FABL.isSmallBiasedSample_zero, FABL.measure_randomSmallBiasFrequencyBadSet_le, FABL.randomSmallBiasFailureSet, FABL.randomSmallBiasFrequencyBadSet, FABL.randomSmallBiasSampleCount, FABL.two_pow_mul_two_exp_neg_four_mul_le_rpow_neg, FABL.uniformPMF}
\difficulty{5}
Let $n$ be a natural number and let $0<\varepsilon<1$. Draw $\lceil
8n/\varepsilon^2\rceil$ points independently and uniformly from the $n$-dimensional cube
over $\mathbb F_2$. The probability that the sample's empirical means of the nontrivial
Walsh characters fail to all have absolute value at most $\varepsilon$ is at most
$2^{-n}$.
\end{theorem}

\begin{theorem}[Bias of a random empirical density]
\lean{FABL.measure_randomSmallBiasDensity_not_isBiased_le}
\label{FABL.measure_randomSmallBiasDensity_not_isBiased_le}
\leanok
\uses{FABL.F₂Cube, FABL.ProbabilityDensity.IsBiased, FABL.measure_randomSmallBiasSample_not_isBiased_le, FABL.randomSmallBiasDensity, FABL.randomSmallBiasDensity_isBiased_iff, FABL.randomSmallBiasFailureSet, FABL.randomSmallBiasSampleCount, FABL.randomSmallBiasSampleCount_pos, FABL.uniformPMF}
\difficulty{3}
Let $n$ be a positive integer and let $0<\varepsilon<1$. Draw $\lceil
8n/\varepsilon^2\rceil$ independent uniform points from the $n$-dimensional cube over
$\mathbb F_2$, and form their empirical probability density relative to the uniform
distribution on that cube. The probability that this density has some nonzero Fourier
coefficient of absolute value greater than $\varepsilon$ is at most $2^{-n}$.
\end{theorem}
